\documentclass{article}

\PassOptionsToPackage{numbers,compress}{natbib}
\usepackage[preprint]{neurips_2024}

\usepackage[utf8]{inputenc}
\usepackage{amsmath}
\usepackage{amssymb}
\usepackage{mathtools}
\usepackage{amsthm}
\usepackage{float}
\usepackage{microtype}
\usepackage{graphicx}
\usepackage{booktabs}
\usepackage{algorithm}
\usepackage{algorithmic}
\usepackage{hyperref}
\usepackage[capitalize,noabbrev]{cleveref}
\RequirePackage{doi}
\usepackage{xcolor,colortbl}

\definecolor{darkblue}{rgb}{0.19,0.19,0.62}
\hypersetup{
  colorlinks=true,
  citecolor=darkblue,
  filecolor=darkblue,
  urlcolor=darkblue
}

\newcommand{\model}{Qwen3-4B-Thinking-2507}
\newcommand{\modelhf}{\texttt{Qwen/Qwen3-4B-Thinking-2507}}
\newcommand{\answerextract}{\mathcal{E}}
\newcommand{\scorefn}{\mathcal{J}}

\title{
CSE 151B/251B Final Project Report\\
{\large Model-Intrinsic Improvement of Qwen3-4B for Mathematical Reasoning}
}

\author{%
  Chuheng Xi, Jeffery Thi, Lin Nie\\
  Department of Computer Science\\
  University of California San Diego\\
  \texttt{cxi@ucsd.edu, jthi@ucsd.edu, y1nie@ucsd.edu}
}

\begin{document}
\maketitle

\begin{abstract}
This project studies how far a small open-weight reasoning model, \model, can be pushed on diverse mathematical reasoning tasks without external tools at inference time. The competition setting is challenging because it mixes high-school, undergraduate, and graduate-level problems, and it requires the same 4B-parameter model to answer both multiple-choice and free-response questions. Our current system improves the starter pipeline through model-intrinsic methods: problem-type-specific prompts, strict final-answer formatting, robust answer extraction, self-consistency voting, and error analysis. On the public labeled data, our first mixed 100-problem prompt experiments reached 44.0\% local accuracy, compared with 41.0\% for the starter-style prompt on the same subset. A later stratified prompt sweep reached 77.0\% on 100 multiple-choice examples with a match-back few-shot prompt and 67.0\% on 100 free-response examples with a multi-answer self-consistency prompt. Our best combined balanced-subset run reached 68/100 = 68.0\% local accuracy, using an MCQ elimination prompt with greedy decoding and a free-response baseline prompt with three-sample voting. The current Kaggle public leaderboard score reported by our team is 60.7\%. These results show that careful alignment between prompt format, answer extraction, and decoding strategy can substantially improve a small reasoning model before any parameter updates. The next stage is to use the collected failure cases to build supervised fine-tuning data and train a LoRA adapter.
\end{abstract}

\section{Introduction}
\label{sec:intro}

\paragraph{Problem definition.}
The task is mathematical answer generation with a fixed language model. Each input is a problem statement $x_i$. Some examples also include a set of answer options $O_i = \{o_{i,A}, o_{i,B}, \ldots\}$, in which case the required output is a single option letter. Other examples are free-response problems whose output may be a number, symbolic expression, equation, interval, ordered list, unordered list, or multiple values corresponding to several \texttt{[ANS]} blanks. The model must generate a textual solution $y_i$, and the grader extracts a final answer $\hat{a}_i = \answerextract(y_i)$ from that text. A prediction is correct when the extracted answer matches the gold answer under the competition-style judger $\scorefn$.

Unlike standard classification, this task combines reasoning, generation, and formatting. The model must understand the problem, choose a method, perform multi-step calculations, and express the final answer in a form that the evaluator can parse. For multiple-choice questions, it is not enough to compute the correct value; the model must map that value back to the option list and output the corresponding letter. For free-response questions, it must decide whether the expected answer should be exact or approximate, how many sub-answers are required, and how to write those answers in a single final box.

\paragraph{Problem significance.}
Reliable mathematical reasoning is important because many real applications require exact quantitative logic rather than fluent explanation alone. A stronger small math model could support tutoring, homework feedback, scientific calculation, data analysis, financial reasoning, and engineering design. In these settings, a superficially plausible answer is often not useful; the final numeric or symbolic result must be correct.

Small models are especially important for deployment. Larger models can be expensive to run and difficult to use on consumer hardware. Our project focuses on \model, a 4B-parameter model, because improving this model would make math reasoning more accessible on limited GPUs. A strong result would show that some reasoning gains come from better prompting, extraction, sampling, and targeted fine-tuning rather than only from scaling model size.

\paragraph{Technical challenge.}
The main technical difficulty is that the benchmark is heterogeneous. The public set contains 1126 labeled problems: 375 multiple-choice problems and 751 free-response problems. The private set contains 943 unlabeled problems: 300 multiple-choice and 643 free-response. The JSON rows only expose \texttt{id}, \texttt{question}, \texttt{options}, and, for public data, \texttt{answer}. They do not provide explicit topic, difficulty, or answer-type labels, so the system must infer the needed behavior from the text itself.

Evaluation is also fragile. Many incorrect local scores are caused not only by wrong reasoning, but also by answer extraction mismatches. We observed cases where the model derived the right relation but omitted one requested blank, rounded an exact answer too aggressively, boxed the numeric value instead of the option letter, or used a notation that was mathematically equivalent but not accepted by the scorer. This makes prompt and post-processing design part of the modeling problem.

Finally, the engineering budget is constrained. The final responses must come from \model, and the team runs experiments on local hardware with 32GB of VRAM. Multi-sample decoding improves robustness but multiplies runtime: generating roughly 1000 private-set answers with three samples per problem can take several hours. This forces us to use small, repeatable validation sweeps before committing to full private-set submissions.

\paragraph{Contributions.}
Our work improves the starter code in five ways.
\begin{itemize}
  \item We built a reproducible vLLM-based inference pipeline around \modelhf, including scripts for public-set prompt experiments, stratified prompt sweeps, best-combination sweeps, and private CSV generation.
  \item We separated multiple-choice and free-response prompting. This matters because MCQ examples require compute-then-match behavior, while free-response examples require exact answer formatting and multi-answer handling.
  \item We implemented prompt families for baseline, strict formatting, verification, concise reasoning, few-shot formatting, option elimination, match-back selection, and multi-answer free-response generation.
  \item We added robust local extraction and scoring logic: MCQ predictions are recovered from boxed letters, option phrases, or boxed option values; free-response predictions are compared through the project judger and normalized before self-consistency voting.
  \item We performed systematic error analysis and used it to guide later prompts. The main observed failures were MCQ letter mismatch, format/extraction errors, overthinking, arithmetic mistakes, rounding mismatch, and missing or misaligned multi-answer outputs.
\end{itemize}

\section{Related Work}
\label{sec:relatedwork}

\paragraph{Chain-of-thought prompting.}
Chain-of-thought prompting encourages language models to write intermediate reasoning steps before the final answer, and it has been shown to improve performance on multi-step reasoning tasks \citep{wei2022chain}. Our system uses this idea, but the experiments show that more reasoning is not always better. Long verification prompts sometimes increased overthinking and truncation, while concise prompts often performed better on the same local subset.

\paragraph{Self-consistency.}
Self-consistency samples several reasoning paths and chooses the most common final answer \citep{wang2023selfconsistency}. This is directly relevant to our project because \model{} can produce different answers under stochastic decoding. We apply self-consistency at the extracted-answer level instead of voting over raw text. For free-response problems, this means normalizing each extracted answer before voting; for MCQ, this means voting over letters.

\paragraph{Math reasoning benchmarks and fine-tuning.}
Datasets such as GSM8K and MATH helped establish mathematical reasoning as a major language-model benchmark \citep{cobbe2021training,hendrycks2021measuring}. These works motivate our focus on both final-answer accuracy and failure analysis. For future parameter-efficient training, we plan to use LoRA, which adapts a frozen base model through low-rank trainable matrices instead of updating all parameters \citep{hu2022lora}. This is a practical fit for our hardware budget.

\section{Methods}
\label{sec:methods}

\paragraph{Problem setting.}
Let $D = \{(x_i, O_i, a_i)\}_{i=1}^{N}$ be the public labeled dataset. Here $x_i$ is the problem statement, $O_i$ is either an option list or empty, and $a_i$ is the gold answer. A prompt template $p$ maps each example into a chat input $T_p(x_i, O_i)$. The fixed base model with parameters $\theta_0$ defines a conditional distribution
\begin{equation}
  y_i \sim \pi_{\theta_0}(\cdot \mid T_p(x_i, O_i)).
\end{equation}
The final prediction is not the full response $y_i$, but the extracted answer $\hat{a}_i = \answerextract(y_i, O_i)$. Local accuracy is
\begin{equation}
  \mathrm{Acc}(D,p) =
  \frac{1}{N}\sum_{i=1}^{N}
  \mathbf{1}\left[\scorefn(\hat{a}_i, a_i, O_i)=1\right].
\end{equation}
For the work completed so far, $\theta_0$ is fixed and we optimize the prompt template, decoding parameters, answer extractor, and voting rule. For the planned LoRA stage, we will train adapter parameters $\phi$ by minimizing the supervised next-token loss
\begin{equation}
  \mathcal{L}_{\mathrm{SFT}}(\phi) =
  -\sum_{i=1}^{M}\sum_{t=1}^{|s_i|}
  \log \pi_{\theta_0+\Delta_{\phi}}
  \left(s_{i,t}\mid T_p(x_i,O_i), s_{i,<t}\right),
\end{equation}
where $s_i$ is a curated solution target with a consistent final-answer format.

\paragraph{Prompt design.}
The starter prompt uses a single general instruction: solve step-by-step and put the final answer in \verb|\boxed{}|. We found this too weak for two reasons. First, MCQ problems often led the model to box the computed value instead of the option letter. Second, multi-answer free-response problems often led the model to answer only the final blank or to split answers across multiple boxes.

We therefore split the prompt space by problem type. For MCQ problems, the best-performing prompt family makes the model solve the problem, explicitly map the computed answer to the option list, and end with a single boxed capital letter. We also tested elimination prompts that ask the model to discard impossible options by sign, magnitude, units, or constraints. For free-response problems, the best prompts emphasize one final boxed answer, exact notation, comma-separated multiple answers, and no labels or units inside the box. The multi-answer prompt explicitly asks the model to count the requested \texttt{[ANS]} blanks before writing the final answer.

\paragraph{Self-consistency and voting.}
For stochastic decoding, we sample $K$ responses for the same prompt:
\begin{equation}
  y_i^{(k)} \sim \pi_{\theta_0}(\cdot \mid T_p(x_i,O_i); \tau, p_{\mathrm{top}}, k_{\mathrm{top}}),
  \quad k=1,\ldots,K.
\end{equation}
Each response is converted into a vote key $v_i^{(k)}$. For MCQ, the key is the extracted option letter. For free response, the key is the normalized extracted answer. The selected response is the first sample whose key equals the majority key. If no answer can be extracted from any sample, we fall back to the first response. This keeps the submitted output as a valid raw model response while making the choice depend on normalized answer agreement.

\begin{algorithm}[H]
\caption{Inference with extraction-level self-consistency}
\label{alg:sc}
\begin{algorithmic}[1]
\STATE Build a chat prompt $T_p(x_i,O_i)$ using an MCQ or free-response template.
\STATE Generate $K$ responses from \model{} with the selected sampling configuration.
\STATE Extract an answer key from each response.
\STATE Normalize keys using letter extraction for MCQ and the math judger for free response.
\STATE Choose the most frequent non-empty key.
\STATE Return the first raw response whose key matches the winning key.
\end{algorithmic}
\end{algorithm}

\paragraph{Answer extraction and error analysis.}
Our extraction logic is designed to match the competition metric more closely than raw string comparison. For MCQ, the extractor first searches for a boxed letter after the reasoning section. If the box contains a value rather than a letter, it tries to match that value against the option list. It then falls back to explicit phrases such as ``option C'' or ``answer is C.'' For free-response problems, we use the project \texttt{Judger} with \texttt{strict\_extract=False}, normalize LaTeX variants such as \verb|\dfrac| and \verb|\tfrac|, and split comma-separated answers when multiple gold values are expected.

We also label failures. The prompt experiment script writes a 20-example wrong-answer sheet for each run and classifies failures as format error, arithmetic error, conceptual error, overthinking, MCQ mismatch, or extraction issue. The sweep script records a simpler taxonomy: correct, truncated, no-answer, out-of-range, wrong, or judge error. These labels are useful because they separate modeling failures from interface failures.

\section{Experiments}
\label{sec:experiments}

\subsection{Baselines}
\label{sec:baselines}

We compare against the starter-style prompt and several inference-time baselines. The starter baseline uses a shared step-by-step instruction and requires a final boxed answer. The prompt-engineered baselines are shown below.
\begin{itemize}
  \item \textbf{Strict format}: requires exactly one boxed final answer and forbids units, labels, or intermediate boxes.
  \item \textbf{Verify}: asks the model to solve, then check the result by substitution, re-derivation, or a sanity check.
  \item \textbf{Concise}: asks for short reasoning to reduce overthinking and truncation.
  \item \textbf{Few-shot}: prepends a small example showing the expected final-answer format.
  \item \textbf{MCQ eliminate}: asks the model to eliminate clearly wrong options before selecting a letter.
  \item \textbf{MCQ match-back}: asks the model to map the computed answer to the option list before boxing the letter.
  \item \textbf{Free multi-answer}: asks the model to count all requested blanks and output all answers in order.
\end{itemize}

\subsection{Evaluation}
\label{sec:evaluation}

The official metric is Kaggle unified accuracy on hidden labels. Locally, we evaluate on the public labeled set using the same high-level principle: extract the final answer and compare it with the gold answer. We use two local evaluation regimes. The first uses the first 100 public examples, which contain 38 MCQ and 62 free-response problems. This was useful for quick prompt iteration. The second uses stratified subsets, typically 100 MCQ and 100 free-response examples or a balanced 50/50 combined subset, which prevents the larger free-response class from hiding MCQ-specific effects.

Because local scoring can differ from Kaggle scoring, we treat local accuracy as a development signal rather than a final claim. The Kaggle leaderboard score is the authoritative estimate of generalization. Our current team-reported public leaderboard score is 60.7\%.

\subsection{Implementation Details}
\label{sec:implementation_details}

All completed experiments use \modelhf{} through vLLM and Hugging Face tokenizers. The initial prompt experiments use bitsandbytes INT8 loading, \texttt{max\_model\_len=16384}, \texttt{max\_tokens=4096}, \texttt{temperature=0.6}, \texttt{top\_p=0.95}, and \texttt{top\_k=20}. The later sweep script uses longer context with \texttt{max\_model\_len=65536}, \texttt{max\_tokens=32768}, prefix caching, and GPU memory utilization 0.85. The best-combination script uses bitsandbytes loading with \texttt{max\_model\_len=16384} and supports both greedy decoding and three-sample self-consistency.

The private-submission script currently defaults to \texttt{match\_back\_few\_shot} for MCQ and \texttt{multi\_answer} for free-response, both with three-sample self-consistency. These defaults reflect the strongest type-specific sweep results, although private-set accuracy cannot be measured locally because the private labels are hidden.

The main decoding configurations are:
\begin{itemize}
  \item \textbf{Greedy}: \texttt{temperature=0.0}, \texttt{top\_p=1.0}, \texttt{top\_k=-1}, one sample.
  \item \textbf{Stochastic single sample}: \texttt{temperature=0.7}, \texttt{top\_p=0.95}, \texttt{top\_k=20}, one sample.
  \item \textbf{Self-consistency}: \texttt{temperature=0.7}, \texttt{top\_p=0.95}, \texttt{top\_k=20}, three samples with majority vote.
\end{itemize}
The experiments were run on a local machine with 32GB of VRAM. Full private-set generation with three samples per problem is expensive; in our current setup, generating roughly 1000 answers can take about 4--5 hours.

\subsection{Results}
\label{sec:results}

\paragraph{Initial mixed prompt experiments.}
\Cref{tab:mixed-prompts} summarizes representative 100-problem runs from \texttt{results/prompt\_experiments/summary.csv}. The strongest initial run was \texttt{split\_concise\_v1} at 44.0\%, a 3 point absolute improvement over a 41.0\% starter run on a comparable 100-problem subset. Verification and detailed reasoning did not reliably help; in several runs they reduced free-response accuracy, likely because they encouraged longer answers and more opportunities for extraction mistakes.

\begin{table}[H]
\caption{Initial mixed 100-problem prompt experiments. Each run uses 38 MCQ and 62 free-response problems from the public set.}
\label{tab:mixed-prompts}
\centering
\begin{tabular}{lrrrrr}
\toprule
Prompt & Overall & MCQ & Free & MCQ correct & Free correct \\
\midrule
Starter & 41.0\% & 44.7\% & 38.7\% & 17/38 & 24/62 \\
Concise & 42.0\% & 47.4\% & 38.7\% & 18/38 & 24/62 \\
Split concise, $\tau=0.3$ & 44.0\% & 50.0\% & 40.3\% & 19/38 & 25/62 \\
Split concise, free $K=4$ & 40.0\% & 34.2\% & 43.5\% & 13/38 & 27/62 \\
Validate random & 26.0\% & 31.6\% & 22.6\% & 12/38 & 14/62 \\
\bottomrule
\end{tabular}
\end{table}

\paragraph{Independent prompt sweep.}
The second experiment separately sweeps MCQ and free-response prompts on stratified 100-example subsets. \Cref{tab:stratified-sweep} shows the best variants. MCQ performance is much higher than free-response performance, suggesting that the model often can solve or guess among options when the final interface is correct. The best MCQ prompt, match-back with a few-shot example, reaches 77/100. The best single-sample free-response prompt reaches 57/100, while the multi-answer self-consistency test reaches 67/100.

\begin{table}[H]
\caption{Best stratified prompt sweep results from \texttt{results/sweep} and \texttt{results/sweep\_sc\_test}. MCQ and free-response rows are evaluated on separate 100-example subsets.}
\label{tab:stratified-sweep}
\centering
\begin{tabular}{llllrrr}
\toprule
Type & Prompt & Decode & Accuracy & Correct & Truncated & Wrong \\
\midrule
MCQ & match\_back\_few\_shot & stochastic $K=1$ & 77.0\% & 77/100 & 4 & 19 \\
MCQ & strict\_format & stochastic $K=1$ & 72.0\% & 72/100 & 11 & 17 \\
MCQ & eliminate & stochastic $K=1$ & 72.0\% & 72/100 & 9 & 19 \\
Free & strict\_few\_shot & stochastic $K=1$ & 57.0\% & 57/100 & 2 & 41 \\
Free & baseline & stochastic $K=1$ & 53.0\% & 53/100 & 4 & 43 \\
Free & multi\_answer & self-consistency $K=3$ & 67.0\% & 67/100 & 2 & 31 \\
\bottomrule
\end{tabular}
\end{table}

\paragraph{Best combined local configuration.}
After the independent sweep, we tested a combined configuration on a balanced 100-problem subset using \texttt{run\_sweep\_best\_combo.py}. The best recorded combination used MCQ elimination with greedy decoding and free-response baseline with three-sample voting. It achieved 68/100 = 68.0\% local accuracy, with 36/50 = 72.0\% on MCQ and 32/50 = 64.0\% on free-response. This is our strongest directly combined local result so far.

\begin{table}[H]
\caption{Best balanced combined local result from \texttt{results/sweep\_best100/summary.csv}.}
\label{tab:best-combo}
\centering
\begin{tabular}{llllr}
\toprule
Subset & Prompt & Decode & Accuracy & Correct \\
\midrule
MCQ & eliminate & greedy $K=1$ & 72.0\% & 36/50 \\
Free & baseline & self-consistency $K=3$ & 64.0\% & 32/50 \\
Combined & MCQ + Free & mixed & 68.0\% & 68/100 \\
\bottomrule
\end{tabular}
\end{table}

\paragraph{Qualitative failure analysis.}
The wrong-answer sheets show several recurring failure modes. First, MCQ prompts can still choose the wrong letter even when the response follows the requested format. For example, one wrong MCQ response computed an answer of 3 and correctly mapped it to option G, but the gold option was E; this is a genuine reasoning or option-interpretation failure, not a formatting failure. Second, free-response models often answer only the final requested blank. In one two-part profit problem, the model produced the correct linear equation $y=470x-390$ in its reasoning but boxed only the second value $\frac{456}{47}$, so the extracted answer missed the first gold value. Third, rounding remains a problem: the model sometimes outputs 7.21 or 7091.67 when the gold answer is a longer decimal or exact expression. Finally, some responses ignore the single-box convention and use \texttt{[ANS]} text or prose after the reasoning, which can confuse extraction.

\section{Discussion}
\label{sec:discussion}

\paragraph{Achievements.}
The main achievement so far is a working inference-time improvement pipeline for \model. The starter-style prompt on the early mixed subset reached 41.0\%, while our best early prompt reached 44.0\%. The stronger stratified sweep showed that task-specific prompting matters more: MCQ match-back few-shot reached 77.0\%, and free-response multi-answer self-consistency reached 67.0\% on their respective 100-example subsets. The best balanced combined local run reached 68.0\%, and our current Kaggle public leaderboard score is 60.7\%. These gains were achieved without updating model weights.

\paragraph{Bottlenecks and mitigation.}
The largest bottleneck is inference cost. Self-consistency improves robustness, but it multiplies generation time because each problem requires multiple full reasoning traces. On our 32GB VRAM setup, a full private-set three-sample run can take 4--5 hours. We mitigate this by first running small public validation sweeps, using stratified subsets, and applying three-sample voting only where it appears useful.

The second bottleneck is evaluator alignment. Some local failures are not pure math failures. They come from missing blanks, wrong answer order, extra labels, decimal-versus-exact mismatch, or MCQ value-versus-letter mismatch. We addressed this by splitting prompts by question type, adding strict final-answer instructions, improving extraction fallback logic, and introducing a multi-answer prompt for free-response questions.

\paragraph{Plans forward.}
The next step is to convert error analysis into training data. We will collect incorrect and borderline public-set examples, label the failure type, and write corrected solution targets with a consistent final-answer format. We will then train a LoRA adapter with supervised fine-tuning so the model learns not only more math patterns, but also the exact output style required by the evaluator. If time permits, we will test simple reinforcement learning or rejection sampling using the local judger as a reward signal, but this is lower priority than a clean SFT dataset and reproducible inference pipeline.

\paragraph{Team responsibilities.}
All team members contributed to strategy discussion, prompt design, and milestone writing. Chuheng Xi focused on running \model{} locally, configuring vLLM on the 32GB VRAM machine, generating submissions, and testing high-cost inference settings. Jeffery Thi focused on organizing experimental ideas, maintaining written notes, and translating observations into report structure. Lin Nie focused on implementation, including multi-sample generation, voting, prompt sweeps, and adapting code to evaluate different prompt families.

\paragraph{Timeline.}
\begin{itemize}
  \item Week 2: Cleaned the baseline setup, unified MCQ and free-response formatting, and established initial local accuracy.
  \item Week 3: Tested temperature, top-$p$, and multi-sample decoding; added majority voting over extracted answers.
  \item Week 4: Added few-shot, strict-format, concise, verify, eliminate, match-back, and multi-answer prompts.
  \item Week 5: Built wrong-answer sheets and classified failures into formatting, arithmetic, conceptual, overthinking, MCQ mismatch, and extraction categories.
  \item Week 6: Ran stratified sweeps for MCQ and free-response prompts and identified the best prompt families.
  \item Week 7: Tested balanced combined configurations and prepared the private submission pipeline.
  \item Week 8: Construct a curated failure dataset for supervised fine-tuning.
  \item Week 9: Train and validate a LoRA adapter with consistent solution and answer formatting.
  \item Week 10: Run final local and Kaggle evaluations, compare inference-only and fine-tuned systems, and finalize the report.
\end{itemize}

\section*{LLM Usage}
We used LLMs as programming and writing assistants. LLM assistance helped brainstorm prompt variants, draft report language, and inspect or refactor experiment scripts. All experimental numbers in this report come from our repository outputs or team submissions, and the authors take responsibility for the final content, claims, and code.

\bibliography{references.bib}
\bibliographystyle{plainnat}

\end{document}
